\section{The paper in brief}\label{sec:01_brief}

\subsection{The question and the answer}\label{sec:01_brief-question}

\textbf{Question.} A takeover target has one bidder that is ready to bid (the \emph{incumbent}). A second
buyer (the \emph{challenger}) must pay a cost to prepare a bid. Does a stronger incumbent keep the
challenger out?

\textbf{Standard answer.} Weakly, yes. A stronger rival lowers what the challenger can expect to earn. Thus entry weakly falls (Fishman 1988; Hirshleifer and Png 1989). This is the
\emph{deterrence effect}.

\textbf{This paper's answer.} Not always. The challenger watches the target's stock price before it
decides. A stronger incumbent can make that price more informative about the challenger's quality class.
If this information effect is large, a stronger incumbent \emph{attracts} the challenger. The title says
this: competition creates competition.

\subsection{The mechanism as a chain}\label{sec:01_brief-mechanism}

Under the class ranking, the auction splits the class value gap between two claims:
\begin{itemize}
\item the \textbf{acquisition claim}: what the challenger keeps if it prepares, bids and wins;
\item the \textbf{traded claim}: what the auction pays to the target's shareholders. The stock trades on
this claim.
\end{itemize}
A stronger incumbent moves these two claims in opposite directions. It weakly lowers the challenger's gross
profit at every belief. It weakly raises the \emph{target-payoff spread} \(\DT=t_H-t_L\), the gap in expected
target proceeds between a high-class and a low-class challenger. An outside investor knows the class and
trades on the spread. The investor shares \emph{one order slot} with liquidity traders, so no order
reveals the class, and the investor keeps a \emph{residual advantage} after every order.
Figure~\ref{fig:01_brief-chain} shows the chain.

\begin{figure}[htbp]
\centering
\begin{tikzpicture}[
  box/.style={draw=intuit, fill=intuit!7, rounded corners=2pt, align=center, text width=3.75cm,
    minimum height=1.45cm, font=\small},
  redbox/.style={draw=caution, fill=caution!7, rounded corners=2pt, align=center, text width=3.75cm,
    minimum height=1.45cm, font=\small},
  arr/.style={-{Stealth[length=2.2mm]}, thick, intuit},
  node distance=0.75cm and 0.55cm]
\node[box] (a) {Stronger incumbent\\(\(F\) falls; \(r\) rises)};
\node[box, right=of a] (b) {Larger spread \(\DT=t_H-t_L\) in target proceeds between the classes};
\node[box, right=of b] (c) {Residual advantage: a correctly signed order earns at least
  \(\rho\,\kappa(\pi)\DT>k\)};
\node[box, below=of c] (d) {Informative orders: class \(H\) buys, class \(L\) sells};
\node[box, left=of d] (e) {A purchase lifts the belief to the top posterior \(M(\pi)\), above
  \(\tau\)};
\node[box, left=of e] (f) {High-cost challenger prepares after a purchase: entry rises above \(\rho\)};
\node[redbox, below=of f] (g) {Direct effect: gross profit \(B_F(\mu)\) falls at every \(\mu\)
  (deterrence)};
\draw[arr] (a) -- (b);
\draw[arr] (b) -- (c);
\draw[arr] (c) -- (d);
\draw[arr] (d) -- (e);
\draw[arr] (e) -- (f);
\draw[-{Stealth[length=2.2mm]}, thick, caution] (a.west) -- ++(-0.35,0) |- (g.west);
\draw[-{Stealth[length=2.2mm]}, thick, caution, dashed] (g.east) -| (e.south)
  node[pos=0.78, right, font=\footnotesize, text=caution, align=left] {raises \(\tau\):\\works against entry};
\end{tikzpicture}
\caption{The two forces. Blue: the information channel. Red: the direct deterrence channel. The bound
\(\rho\,\kappa(\pi)\DT\) applies to the gross return of a correctly signed nonzero order (a purchase in
class \(H\), a sale in class \(L\)). The paper finds an open set of primitives where blue wins.}
\label{fig:01_brief-chain}
\end{figure}

\subsection{The main result in numbers}\label{sec:01_brief-numbers}

The benchmark has \(h=\reg{prim_h}\), \(\ell=\reg{prim_ell}\), \(\pi=\reg{prim_pi}\), \(p=\reg{prim_p}\),
\(\rho=\reg{prim_rho}\), \(c_H=\reg{prim_cH}\) and \(k=\reg{prim_k}\). The incumbent is uniform on
\([0,r]\). The slot has \(1+u=\eu\) and \(\beta=\reg{beta_B}\), so the presence probability is
\(\lambda=\reg{lambda_B}\) and the top posterior is \(M(1/2)=\reg{M_half_B}\). Table~\ref{tab:01_brief-numbers}
compares the three economies of Proposition~2.

\begin{table}[htbp]
\centering\small
\begin{tabular}{@{}lccc@{}}
\toprule
 & Weak (\(r_0=\reg{prim_r0}\)) & Strong (\(r_1=\reg{prim_r1}\)) & Very strong (\(r_2=\reg{prim_r2}\))\\
\midrule
Gross profit at the prior, \(B_r(1/2)\) & \reg{Bprior_r0} & \reg{Bprior_r1_B} (lower) & \reg{Bprior_r2}\\
Target-payoff spread, \(\DT\) & \reg{Delta_r0} & \reg{Delta_r1} (\(40\times\)) & \reg{Delta_r2}\\
Preparation threshold, \(\tau(r)\) & \reg{tau_r0} & \reg{tau_r1} & \reg{tau_r2}\\
Top posterior, \(M(1/2)\) & \reg{M_half_B} & \reg{M_half_B} & \reg{M_half_B}\\
Investor orders, classes \(H\) and \(L\) & \((0,0)\) & \((1,-1)\) & \((1,-1)\)\\
Price informative? & no & yes & yes\\
High-cost preparation & none & after a purchase & none\\
Entry \(\Ent\) & \reg{entry_r0} & \textbf{\reg{entry_r1_B}} & \reg{entry_r2}\\
High-class ownership \(\Own_H\) & \reg{own_r0} & \textbf{\reg{own_r1_B}} & \reg{own_r2_B}\\
\bottomrule
\end{tabular}
\caption{The main result at the benchmark \paper{Tables~1 and~3}.}
\label{tab:01_brief-numbers}
\end{table}

From \(r_0\) to \(r_1\), gross profit falls by about \(10.7\) percent, while the spread rises \(40\) times. % verify/03_auction.py Entry more than doubles, by a factor of
\(2.09\). % verify/01_brief.py
At \(r_2\) the price stays informative, but \(\tau(r_2)=\reg{tau_r2}\) lies just above
\(M(1/2)=\reg{M_half_B}\). So even a purchase cannot recruit a high-cost challenger, and entry falls back
to \(\rho\). Entry is \emph{non-monotone} in incumbent strength.

\begin{intuition}[why entry doubles while profit falls]
Profit falls at \emph{every} belief, but the belief that matters changes. At \(r_0\) the price says nothing,
so the challenger decides at the prior \(1/2\), and \(B_{r_0}(1/2)=\reg{Bprior_r0}<c_H=6\). At \(r_1\) a
purchase moves the belief to \(\reg{M_half_B}\), and \(B_{r_1}(M)\) exceeds \(6\) by
\(\reg{margin_recruit_B}\). A purchase occurs with probability \(b+\lambda/2\) at the prior \(1/2\), so entry is
\(\rho+(1-\rho)(b+\lambda/2)=\reg{entry_r1_B}\).
\end{intuition}

\subsection{Seven claims that you must be able to defend}\label{sec:01_brief-claims}

Use this list to locate each proof.
\begin{enumerate}
\item \textbf{Two returns to information} \paper{Prop.~1}.
A stronger incumbent weakly lowers each class's gross profit.
Under the class ranking, it weakly raises the target-payoff spread.
The proof uses payment integrals and Tonelli's theorem.
\Cref{sec:03_auction-prop1,sec:03_auction-prop1c} state the exact scope.
\item \textbf{Posterior bounds} \paper{Lemma~A.1; Cor.~A.1}.
Every posterior lies in \([m(\pi),M(\pi)]\), under every investor profile, including mixed profiles and wrong-signed orders.
The zero order needs \(\beta\le1/3\).
With \(g_H\ge g_L\), \(B(M(\pi))<c_H\) excludes high-cost preparation under every profile
(\Cref{sec:04_slot_beliefs}).
\item \textbf{Price sufficiency} \paper{Lemma~A.2}.
If \(\DT>0\), equal prices imply equal posteriors because entry is at least \(\rho>0\).
If also \(g_H>g_L\), prices rise strictly with the posterior
(\Cref{sec:04_slot_beliefs-sufficiency}).
\item \textbf{Global trading bounds} \paper{Lemma~A.3}.
Under every candidate, \(\rho\kappa(\pi)\DT\le K_H,K_L\le\DT\).
These are gross returns of correctly signed nonzero orders.
Zero earns zero; a wrong-signed order earns at most \(-k\) net.
The trading-cost window (H2) therefore forces no trade at \(F_0\) and full orders at \(F_1\)
(\Cref{sec:05_investor,sec:06_main_result}).
\item \textbf{Information is the cause} \paper{Prop.~A.3}.
At fixed information, a stronger incumbent never raises entry.
The theorem needs \(\mathcal L_i(\theta,R\mid j,S)=G_j\otimes F_i\) and a cost independent of all these variables.
Frozen full orders give entry \reg{p4_frozen_r0_entry} at both \(r_0\) and \(r_1\).
A hidden price gives entry \reg{p4_hidden_r0_entry} at both strengths
(\Cref{sec:07_controls}).
\item \textbf{Coexistence} \paper{Prop.~3}.
Under rules IM and C1, no informed trade and an informative equilibrium coexist exactly on \([r_A,r_N]\).
At the benchmark, \(r_A\approx\reg{rA}\) and \(r_N\approx\reg{rN}\).
Entry on the informative branch rises to \(r_F\), then stays constant until the entry ceiling.
The proof uses algebra and two exact Sturm counts
(\Cref{sec:08_coexistence}).
\item \textbf{Research} \paper{Prop.~4}.
The investor chooses presence at cost \((a/2)\lambda^2\).
On the interior research branch, research, entry and high-class ownership rise with incumbent strength.
At \(r=\reg{str_r160}\), research is \reg{res_lambda_r160} and entry is \reg{res_entry_r160}.
Under C1, this branch coexists with no research
(\Cref{sec:09_research}).
\end{enumerate}

\subsection{A 60-second version for the panel}\label{sec:01_brief-sixty}

\begin{panelnote}
``A stronger rival lowers the challenger's expected auction profit.
I add a stock market before the challenger commits.
Under the class ranking, stronger competition also widens the gap in target proceeds between the two classes.
An informed investor trades on that gap.
Liquidity orders keep each order from fully revealing the class.
When the residual advantage covers the trading cost, informed trade starts.
A purchase can then recruit a high-cost challenger that stays out at its prior.
On an open set of primitives, entry rises with incumbent strength.
The weak economy has one uninformative outcome; the strong economy has one informative outcome with more entry.
At fixed information, entry cannot rise.
The change in information therefore explains the reversal.''
\end{panelnote}

\subsection{How to read this guide}\label{sec:01_brief-read}

\Cref{sec:02_model} sets out the model. \Cref{sec:03_auction} derives every auction payoff.
\Cref{sec:04_slot_beliefs} derives what one order can reveal. \Cref{sec:05_investor} solves
the investor's problem. \Cref{sec:06_main_result} proves the main result step by step.
\Cref{sec:07_controls,sec:08_coexistence,sec:09_research,sec:10_robustness,sec:11_welfare,sec:12_empirical} cover the controls, coexistence, research,
robustness, welfare and sale terms, and empirical implications.

Work in this order: the auction algebra first, because every later step uses \(\DT\) and \(B(\mu)\); then
the posterior bounds; then the trading bounds. With these three tools, the proof of Proposition~2 is a
short sequence of inequalities.

\begin{intuition}[What changed from the Laplace version]
One order slot replaces additive Laplace flow.
Orders are \(\{-1,0,1\}\), and the low preparation cost is zero.
The auction now allows value classes; the point masses remain the benchmark.
The guide adds the full equilibrium path and costly research.
\Cref{sec:10_robustness-sizes} derives the exact comparison of posterior bounds.
\end{intuition}
